\documentclass[11pt]{article}
\usepackage[a4paper,margin=27mm]{geometry}
\usepackage{amsmath,amssymb,mathtools}
\usepackage{fontspec,unicode-math}
\setmainfont{Latin Modern Roman}
\setmathfont{Latin Modern Math}
\usepackage{xurl,hyperref}
\hypersetup{colorlinks=true,urlcolor=blue,linkcolor=blue}
\newcommand{\sref}[2]{\hyperlink{source:#1:#2}{[S#2]}}
\newcommand{\cataloguescope}[1]{\paragraph{Mathematical question.} #1}
\begin{document}
\section{Conjecture 3.23}\label{statement}

\subsection{Definitions and mathematical statement}

\paragraph{Shared notation.}$\mathbb N=\{1,2,\ldots\}$, and $\mathbb Z,\mathbb Q,\mathbb R,\mathbb C$ denote the integers, rationals, reals and complex numbers. Any other conventions are specified locally.


Let $\Lambda\subset\mathbb C$ be a rank-two lattice, with Weierstrass functions $\wp_\Lambda,\zeta_\Lambda$ normalized by $\zeta_\Lambda^{\prime}=-\wp_\Lambda$ and $\zeta_\Lambda(z)=1/z+O(z^3)$ at zero. Its invariants satisfy $(\wp_\Lambda^{\prime})^2=4\wp_\Lambda^3-g_2\wp_\Lambda-g_3$ and $g_2^3-27g_3^2\ne0$. For $0\ne\omega\in\Lambda$, define the constant $\eta=\zeta_\Lambda(z+\omega)-\zeta_\Lambda(z)$. Transcendence degree counts the largest algebraically independent subfamily over $\mathbb Q$.
\paragraph{Statement recorded in the supplied source.}
For every such lattice, nonzero period $\omega$, and $u\in\mathbb C\setminus\Lambda$, is \[\operatorname{trdeg}_{\mathbb Q}\mathbb Q\bigl(g_2,g_3,\pi/\omega,\wp_\Lambda(u),\zeta_\Lambda(u)-(\eta/\omega)u\bigr)\ge2?\] \sref{statement}{2}
\subsection{Short English statement}
Must this field of elliptic invariants, one normalized period and two elliptic-function values have at least two algebraically independent elements?
\subsection{Sources}
\begin{description}
\item[\hypertarget{source:statement:1}{S1}] ULAM.AI and the UnsolvedMath Contributors, UnsolvedMath: A Curated Collection of Open Mathematics Problems, catalogue entry. CC BY 4.0. \url{https://huggingface.co/datasets/ulamai/UnsolvedMath}.
\item[\hypertarget{source:statement:2}{S2}] Michel Waldschmidt, Open Diophantine Problems, arXiv:math/0312440v2 (24 January 2004), Conjecture 3.23. \url{https://arxiv.org/abs/math/0312440}.
\end{description}
\label{end:statement}
\end{document}
